\documentclass[conference]{IEEEtran}
\IEEEoverridecommandlockouts

\usepackage{cite}
\usepackage{amsmath,amssymb,amsfonts}
\usepackage{graphicx}
\usepackage{booktabs}
\usepackage{hyperref}

% \anontrue hides authors and uses the anonymized code link. IEEE ICSC 2027
% does not state its review model, so we submit anonymized; switch to
% \anonfalse for the camera-ready version after acceptance.
\newif\ifanon
\anontrue

\begin{document}

\title{Evaluating LLM-Augmented ML Pipelines: Performance, Stability, and Cost}

\ifanon
\author{\IEEEauthorblockN{Anonymous Author(s)}}
\else
% TODO(authors): replace with the real author list and affiliations.
\author{
  \IEEEauthorblockN{AUTHOR NAMES}
  \IEEEauthorblockA{AFFILIATION \\ City, Country \\ EMAIL}
}
\fi

\maketitle

\begin{abstract}
Large language models are increasingly used as item encoders inside
existing ranking pipelines, often on the assumption that a bigger
encoder gives better features. We test that assumption under controlled
conditions. Six frozen encoders feed the same personalized
learning-to-rank pipeline: TF-IDF, a 22M-parameter sentence encoder, a
335M-parameter embedding encoder, a 3.8B chat LLM used through its
mean-pooled hidden states, and a 4B LLM trained for embedding, plus a
hybrid. We evaluate them on MovieLens, Amazon Reviews, and Yelp (up to
10{,}537 users) with popularity and matrix-factorization baselines,
uniform and popularity-matched negatives, random and temporal splits,
two scorers, paired per-user significance and equivalence tests, and a
like-for-like latency benchmark. Three findings emerge. First,
evaluation design can erase the question: a scorer that sees only item
features reduces to a popularity model, and easy negatives let
popularity match every encoder. Second, with a personalized scorer and
hard negatives, moving from TF-IDF to the small sentence encoder is the
most consistent gain, significant on all three datasets (on Amazon only
under one of the two scorers).
Above it, gains are small and need a large sample to detect: on Yelp
the 4B embedding model adds 0.006--0.016 NDCG@10 at roughly
45--70$\times$ the per-item latency, and on MovieLens and Amazon no gain
is detectable. Third, a chat LLM's raw hidden states are unreliable.
They are nearly parallel on every dataset, significantly worse than
TF-IDF on short metadata, and competitive on long reviews only under a
tree-based ranker. An LLM of the same size trained for embedding has no
such failure. We turn these results into a cost-aware checklist for
adding LLM components to ranking pipelines.
\end{abstract}

\begin{IEEEkeywords}
large language models, recommender systems, learning to rank, embeddings, quality-cost tradeoff, empirical evaluation
\end{IEEEkeywords}

\section{Introduction}
\label{sec:introduction}

Large language models (LLMs) are being incorporated into machine learning
systems at a rapid pace: as embedding providers for retrieval and ranking,
as feature-enrichment components that turn raw text into structured
signals, and increasingly as reasoning components sitting directly in a
production decision pipeline \cite{llm4rec-survey}. Recommender-systems
research has a documented history of gains that shrink or vanish under
controlled comparison: reproducibility studies found that many reported
neural improvements did not survive against properly tuned simple
baselines \cite{dacrema-progress, rendle-ncf-mf}, that conclusions shift
with splitting, sampling, and tuning choices \cite{sun-rigorous}, and that
sampled-negative metrics can mis-rank methods \cite{krichene-sampled}.
Cost is rarely part of the comparison at all. LLM components inherit
all of these risks and add a new one: a per-item inference cost tens of
times that of the component they replace.

In practice, the behavior of LLM-augmented pipelines is far less
consistent than the enthusiasm around them would suggest. The same
architecture that provides a clear win on one dataset can be a wash, or
even a regression, on another. Practitioners deciding whether to add an
LLM component to an existing ML pipeline need more than a single
accuracy number: they need to know whether the gain is stable across
random seeds and data splits, whether it survives distribution shift, and
whether it is worth the additional latency and cost relative to a much
cheaper classical or small-model baseline.

This paper studies LLMs as \emph{components inside} ML pipelines rather
than as end-to-end systems, and asks a narrower, more actionable question
than ``do LLMs help'': \emph{when} do they help, by how much, how
reliably, and at what cost, relative to the alternatives a practitioner
would otherwise reach for.

\subsection{Contributions}
\begin{itemize}
  \item A controlled comparison of six frozen item encoders spanning
    classical TF-IDF, a 22M-parameter sentence encoder, a 335M-parameter
    embedding encoder, a 3.8B chat LLM used through its raw hidden
    states, and a 4B LLM trained for embedding. All of them feed the
    same personalized learning-to-rank pipeline, under two scorers, two
    negative-sampling regimes, a random and a temporal split, and
    popularity and matrix-factorization reference baselines
    (\S\ref{sec:pipeline-design}--\S\ref{sec:results}).
  \item Paired per-user significance and equivalence tests with a
    pre-stated margin, so that ``no gain'' claims are distinguished from
    ``no evidence'' (\S\ref{sec:evaluation-metrics}).
  \item Three findings, replicated across three datasets. (i) A scorer
    that sees only item features reduces to a popularity model, under
    which every encoder looks alike. (ii) Once ranking is personalized,
    the step from TF-IDF to a small dense encoder is the most
    consistent gain. Larger encoders add a small gain only on long
    review text with $\sim$10K users, at tens of times the cost
    (\S\ref{sec:results}). (iii) A chat LLM's raw hidden states can be
    significantly \emph{worse} than TF-IDF, depending on text length
    and scorer, while an LLM of the same size trained for embedding is
    reliable (\S\ref{sec:failure-analysis}). Within one dataset, the
    advantage of dense over lexical encoders also tracks text redundancy
    (\S\ref{sec:redundancy-analysis}).
  \item A like-for-like latency benchmark, with the same batching, inputs,
    and hardware for every encoder, and a cost-aware checklist for
    practitioners (\S\ref{sec:quality-cost},
    \S\ref{sec:practical-guidelines}).
\end{itemize}

\textbf{A cautionary note on evaluation design.} Our first finding is
methodological: feeding item embeddings to a ranker that sees nothing
about the user makes every encoder look alike, and sampling easy
negatives hides the gap between lexical and dense encoders. Either one
alone would have led us to a wrong conclusion
(\S\ref{sec:movielens-results}).

\textbf{Reproducibility.} Code, experiment scripts, per-user results,
and the scripts that generate every table are available at
\ifanon
\url{https://anonymous.4open.science/r/llm-augmented-ml-pipelines-FC52/}.
\else
\url{CAMERA-READY-CODE-LINK}.
\fi

\section{Related Work}
\label{sec:related-work}

\textbf{Text and LLM representations in recommendation.} Surveys
distinguish LLMs used directly as recommenders from LLMs that supply
embeddings or generated features to a conventional model
\cite{llm4rec-survey}; we study the latter. Text-based item
representations from pretrained encoders have been used to make
sequential recommenders transferable across domains (UniSRec
\cite{unisrec}, Recformer \cite{recformer}), and Yuan et
al.~\cite{yuan-id-vs-modality} show that modality-based item encoders
can match ID-based ones when the encoder is trained end-to-end. Harte et
al.~\cite{harte-llmseqrec} report gains from initializing sequential
recommenders with LLM embeddings, and BLaIR \cite{amazon-reviews-2023}
pretrains sentence encoders on Amazon review text specifically to
connect language and items. These works change the recommender
architecture, fine-tune the encoder, or both. We ask a narrower,
deployment-side question that they leave open: holding a standard
feature-based ranker fixed and using \emph{frozen, off-the-shelf}
encoders, does moving up the capacity ladder (classical
$\rightarrow$ small encoder $\rightarrow$ large encoder $\rightarrow$
LLM) pay for its cost?

\textbf{Decoder LLMs as embedders.} Raw hidden states of decoder-only
LLMs are a weak embedding without adaptation: SGPT \cite{sgpt},
LLM2Vec \cite{llm2vec}, and E5-Mistral \cite{e5-mistral} all add
contrastive training (and, for LLM2Vec, bidirectional attention) to
reach state-of-the-art results on MTEB \cite{mteb}. Our phi4-mini arm
deliberately uses the training-free variant, which is what a
practitioner gets by pointing an existing chat model at an embedding
endpoint; we therefore read its results as a statement about that
practice, not about the ceiling of LLM-derived embeddings
(\S\ref{sec:discussion}).

\textbf{Baselines and evaluation rigor in recommendation.} Several
reproducibility studies found that reported neural gains often vanish
against properly tuned simple baselines such as item-kNN or matrix
factorization \cite{dacrema-progress, rendle-ncf-mf}, and that results
are sensitive to splitting, sampling, and tuning choices that papers
rarely control \cite{sun-rigorous}. Krichene and Rendle
\cite{krichene-sampled} show that metrics computed over sampled
negatives can be inconsistent with full-ranking metrics and are most
lenient when negatives are easy. We respond to these findings directly:
we include popularity and PureSVD \cite{puresvd} reference baselines,
evaluate under both uniform and popularity-matched negatives, and base
every ``no difference'' claim on paired equivalence tests (TOST
\cite{tost}) with a pre-stated margin rather than on overlapping
confidence intervals.

\textbf{Gradient-boosted rankers.} LightGBM \cite{lightgbm} represents
the gradient-boosted tree family widely used as the final scorer in
production ranking stacks; it consumes dense item embeddings and
handcrafted user--item features side by side, which is exactly the
integration point at which a practitioner would add an LLM component.

\section{Problem Setup}
\label{sec:problem-setup}

We study pipelines of the general form
\[
\text{item text} \xrightarrow{\;\text{encoder}\;} \text{item vector}
\;\rightarrow\; \text{user--item features} \;\rightarrow\; \text{scorer}
\;\rightarrow\; \text{ranked list},
\]
in which a frozen text encoder turns each item's text into a vector, the
vector feeds a set of user--item features, and a learned scorer ranks
each user's candidate items. This is the integration point at which a
practitioner typically adds an LLM: the scorer, the candidate lists, and
the training data stay as they are, and only the encoder changes. Holding
everything except the encoder fixed lets us attribute differences in
ranking quality to the item representation alone.

\textbf{What the scorer sees.} For a user $u$ and candidate item $i$ with
encoder vector $\mathbf{e}_i$, the scorer receives
\begin{enumerate}
  \item the item vector $\mathbf{e}_i$ itself;
  \item two \emph{affinity} features: the mean and the maximum cosine
    similarity between $\mathbf{e}_i$ and the vectors of $u$'s training
    positives $P_u$. When $i \in P_u$ (a training row), $i$ is removed
    from $P_u$ before computing these features (leave-one-out), so the
    scorer never trains on a self-similarity it could not see at test
    time;
  \item two \emph{popularity} features computed from training
    interactions only: $\log(1+n_i)$, where $n_i$ is the number of
    training ratings of $i$, and $i$'s mean training rating.
\end{enumerate}
The affinity features are where representation quality can matter per
user: they are computed in each encoder's own embedding space. The
popularity features are identical across encoders and so shift every
pipeline equally. We also report two ablations: the scorer without the
popularity features, and the scorer with the item vector alone. The
latter is a common shortcut when embeddings are dropped into an
existing ranker, and in it the scorer can only learn a single global
score per item
(\S\ref{sec:results}).

\textbf{Scope.} We study one form of LLM augmentation, namely item
embeddings from a frozen encoder. We do not study LLM-generated features
or LLMs used as end-to-end recommenders. All datasets are public
benchmarks standing in for production traffic, so absolute numbers
should be read as comparisons between pipelines on a fixed dataset.

\section{ML Pipeline Design}
\label{sec:pipeline-design}

All pipelines share candidate generation, splits, the scorer and its
hyperparameters, and the feature construction of
\S\ref{sec:problem-setup}. Only the encoder that produces
$\mathbf{e}_i$ differs.

\subsection{Encoders: a capacity ladder}
\label{sec:baseline-pipelines}
\label{sec:large-embedding-pipeline}
\label{sec:llm-augmented-pipelines}
\label{sec:hybrid-pipelines}
We compare six item representations in increasing order of capacity and
cost:
\begin{itemize}
  \item \textbf{TF-IDF}: a 2{,}000-term TF-IDF vector over the item
    text, which is classical and involves no model inference.
  \item \textbf{SBERT}: \texttt{all-MiniLM-L6-v2} \cite{sbert}, a
    22M-parameter sentence encoder. It is the cheapest dense encoder
    and the reference point for every cost comparison.
  \item \textbf{mxbai}: \texttt{mxbai-embed-large}, a 335M-parameter
    dedicated embedding encoder (BERT-style, not an LLM), served
    through Ollama.
  \item \textbf{phi4-mini (raw)}: the 3.8B-parameter decoder-only chat
    LLM \texttt{phi4-mini}, used as an embedder by mean-pooling its
    last-layer hidden states with no embedding-specific training
    (\texttt{llama-server --pooling mean}). This is what a practitioner
    gets by pointing an existing chat model at an embedding endpoint.
  \item \textbf{Qwen3-Embedding-4B (trained)}: a 4B-parameter
    decoder-only LLM \emph{contrastively trained for embedding}, served
    with last-token pooling as its model card specifies. Its size and
    quantization (Q4\_K\_M) match phi4-mini's, so comparing the two
    separates the effect of LLM capacity from the effect of using
    untrained hidden states as an embedding
    \cite{sgpt, llm2vec, e5-mistral}.
  \item \textbf{Hybrid}: TF-IDF concatenated with phi4-mini. The two
    vectors each contribute their own affinity features.
\end{itemize}
SBERT and mxbai accept at most 256 and 512 input tokens respectively,
while review-based item text runs longer. We therefore split long text
into word chunks that fit each encoder (150 and 300 words), embed each
chunk, and mean-pool the chunk vectors, so that every encoder sees the
whole item text. Both LLMs run with a 4{,}096-token context, which
covers every item without chunking. All models run locally on one Apple
M4 machine (16\,GB), and all GGUF models use the same Q4\_K\_M
quantization.

\subsection{Reference baselines}
Two baselines use no item text at all and give the absolute numbers a
reference point. \textbf{Popularity} ranks candidates by $\log(1+n_i)$.
\textbf{PureSVD} \cite{puresvd} factorizes the binarized training
user--item matrix (rank 50, fixed in advance rather than tuned on test
lists) and scores $(u,i)$ by the reconstructed entry. A content pipeline
that cannot beat popularity is not using its item representation for
personalization. A pipeline that cannot beat PureSVD gains nothing from
text over a standard collaborative model.

\subsection{Scorers}
We use two scorers so that conclusions are not tied to one model's
inductive biases. The first is LightGBM's \texttt{LGBMRanker}
\cite{lightgbm} (lambdarank objective, grouped by user, 200 trees, 64
leaves, learning rate 0.05). The second is a pointwise multi-layer
perceptron (scikit-learn \texttt{MLPRegressor}, hidden layers 64 and 32,
standardized inputs, early stopping) that regresses the graded
relevance label. Hyperparameters are fixed across all pipelines and
datasets.

\section{Experimental Setup}
\label{sec:experimental-setup}

\subsection{Datasets}
\label{sec:datasets}

\textbf{MovieLens.} \texttt{ml-latest-small} \cite{movielens}: 100{,}836
ratings (0.5--5.0 stars, timestamped) from 610 users on 9{,}724 movies.
Item text is the title and genre tags (e.g.\ ``Toy Story (1995) Adventure
Animation Children Comedy Fantasy''), which is short, curated metadata.

\textbf{Amazon Reviews.} The \texttt{All\_Beauty} category of Amazon
Reviews 2023 \cite{amazon-reviews-2023}, restricted to its standard
5-core subset: 253 users, 356 items, and 2{,}535 ratings (1--5 stars).
Each item's text is its product title, its category, and up to 10 of its
reviews (title and body, each truncated to 500 characters).

\textbf{Yelp.} Businesses in Boise, ID, from the Yelp Open Dataset, with
at least 5 reviews: 2{,}937 businesses and 105{,}366 ratings (1--5
stars). 10{,}537 users have at least two positives and are evaluated
(median 3 positives). We chose a single mid-sized city to keep the
LLM embedding pass tractable on one machine. Item text is the business
name, its categories, and up to 10 reviews (500 characters each).

\textbf{Leak-free review text.} Review text is strongly correlated with
the star rating that defines relevance. If a held-out user's review were
part of an item's text, the label would leak into the features. We
therefore build item text only from reviews written by users who are
\emph{never evaluated}: reviewers outside the 5-core user set on Amazon,
and users with fewer than two positives on Yelp. No review whose rating
is a test label, under either split, can enter any item's text. After
this restriction the mean item text is 2{,}083 characters on Amazon and
2{,}148 on Yelp. Items with no qualifying review (18 of 356 on Amazon,
99 of 2{,}937 on Yelp) fall back to title and category.

\textbf{Within-dataset redundancy variants.} To test the redundancy
diagnostic (\S\ref{sec:redundancy-analysis}) without changing users, items, or
labels, we also run MovieLens with title-only and genres-only item
text.

\subsection{Task and candidate generation}
\label{sec:candidate-generation}
A rating $\geq 4.0$ is a positive. For each user with at least two
positives, we hold out 20\% of their positives (at least one) for test.
Each training positive is paired with 4 sampled negatives, and each test
list adds 50 sampled negatives drawn from items the user never rated.
Relevance is the rating for positives and 0 for negatives.

\textbf{Two negative-sampling regimes.} Sampled-negative metrics are
most lenient when negatives are easy \cite{krichene-sampled}.
Uniformly sampled unrated items are mostly obscure, so popularity alone
separates them from positives. We therefore report two regimes:
\emph{uniform} negatives, and \emph{popularity-matched} negatives drawn
with probability proportional to each item's rating count, which
removes most of the popularity shortcut. We treat popularity-matched
negatives as the primary setting.

\textbf{Splits.} The \emph{random} split shuffles each user's positives
before holding out 20\%. The \emph{temporal} split holds out each
user's most recent 20\%, so that the model must extrapolate to future
preferences. Popularity features are computed from every rating except
the held-out test positives. Under the temporal split they can
therefore include other users' later ratings, which is a mild,
pipeline-independent form of look-ahead that we note but do not remove.

\subsection{Metrics and statistics}
\label{sec:evaluation-metrics}
For each user we compute NDCG@10 with gain $2^{r}-1$ \cite{ndcg} and
Recall@10. Every configuration runs 5 seeds, and each seed redraws the
split and the negatives. Tables report the mean over seeds of the
per-user average, with a 95\% CI across seeds.

\textbf{Paired tests.} All pipelines see identical users, splits, and
candidate lists in each seed, so every comparison is paired. For each
pipeline pair we average each user's NDCG@10 over seeds and test the
per-user differences with (i) a paired $t$-test for any difference and
(ii) TOST equivalence \cite{tost} against a margin of $\pm0.01$
NDCG@10. The margin was fixed before the main experiments were run; it
equals half of the largest encoder gap observed in pilot runs with the
item-only scorer. Both families are Holm-corrected within each
scorer and configuration. We call a pair \emph{different} if only the
$t$-test rejects, \emph{equivalent} if only TOST rejects,
\emph{negligible} if both reject (a detectable difference inside the
margin), and \emph{inconclusive} if neither rejects.

\textbf{Latency.} Embedding latency is measured in a separate
benchmark, not inside the experiment runs. All four neural encoders
embed the same 256-item sample, with full-text coverage (chunking where
needed), at batch sizes 1 and 32, on the same machine with nothing else
running. Each model gets one warmup call, and three rounds rotate the
model order. We report the median per-item latency over rounds.
SBERT runs on the Metal (MPS) backend, mxbai through Ollama, and the
two GGUF models through llama.cpp servers with four parallel slots,
so a batch request is spread across slots. Every
pipeline runs locally, so the marginal API cost is \$0 and latency is
the cost signal.

\section{Results}
\label{sec:results}

\begin{table*}[t]
\centering
\caption{MovieLens NDCG@10 (LGBMRanker, mean $\pm$ 95\% CI over 5 seeds).
Columns add scorer features left to right under popularity-matched
negatives; the last column uses uniform negatives. Popularity and PureSVD
use no item text.}
\label{tab:main-results}
\begin{tabular}{lcccc}
\toprule
 & \multicolumn{3}{c}{Popularity-matched negatives} & \\
\cmidrule(lr){2-4}
Pipeline & Item only & +Affinity & +Aff.+Pop. & Uniform neg. \\
\midrule
Popularity & 0.379$\pm$0.002 & 0.379$\pm$0.002 & 0.379$\pm$0.002 & 0.731$\pm$0.006 \\
PureSVD & 0.572$\pm$0.008 & 0.572$\pm$0.008 & 0.572$\pm$0.008 & 0.757$\pm$0.005 \\
\midrule
TF-IDF & 0.383$\pm$0.004 & 0.573$\pm$0.007 & 0.576$\pm$0.008 & 0.784$\pm$0.005 \\
SBERT (22M) & 0.380$\pm$0.003 & 0.590$\pm$0.007 & 0.589$\pm$0.008 & 0.785$\pm$0.005 \\
mxbai (335M) & 0.379$\pm$0.004 & 0.586$\pm$0.005 & 0.583$\pm$0.005 & 0.785$\pm$0.005 \\
phi4-mini raw (3.8B) & 0.379$\pm$0.005 & 0.486$\pm$0.006 & 0.486$\pm$0.006 & 0.747$\pm$0.006 \\
Qwen3-Emb.\ (4B) & 0.378$\pm$0.004 & 0.593$\pm$0.007 & 0.587$\pm$0.005 & 0.785$\pm$0.005 \\
Hybrid & 0.380$\pm$0.004 & 0.579$\pm$0.007 & 0.579$\pm$0.009 & 0.783$\pm$0.005 \\
\bottomrule
\end{tabular}

\end{table*}

\subsection{An item-only scorer is a popularity model}
\label{sec:movielens-results}
The ``Item only'' column of Table~\ref{tab:main-results} is the common
shortcut of feeding item embeddings straight into a ranker: the scorer
sees item vectors and nothing about the user. Under popularity-matched
negatives every encoder lands within 0.004 NDCG@10 of plain popularity
(0.378--0.383 vs.\ 0.379), and all are 0.19 below PureSVD. A scorer with no user
input can only learn one global score per item, so it ranks every
user's candidates identically and the item representation is reduced to
a predictor of popularity. This explains why every encoder looked alike
in that setting, and a study run only in this configuration would
conclude, wrongly, that encoder choice does not matter. On Amazon and
Yelp the item-only scorer falls \emph{below} popularity
(Table~\ref{tab:amazon-results}). Once the scorer receives the user--item affinity
features (``+Affinity''), TF-IDF, SBERT, mxbai, and Qwen3 gain 0.19--0.22
(phi4-mini only 0.11), and the encoders separate. Adding the popularity features changes no pipeline
by more than 0.006. Under uniform negatives, popularity alone already
reaches 0.731 on MovieLens, and on Yelp it ties TF-IDF (0.465) with every
encoder within 0.007 of it. This shows how much of an easy-negative
evaluation a non-personalized model can absorb \cite{krichene-sampled}.

\subsection{The capacity ladder on MovieLens}
\label{sec:cross-dataset-synthesis}
\begin{table}[t]
\centering
\caption{Paired per-user tests against SBERT, all datasets and scorers
(personalized scorer, popularity-matched negatives, 5 seeds). Cells are
NDCG@10 of the pipeline minus SBERT. $^{*}$: Holm-corrected $p<0.05$;
$^{\equiv}$: TOST-equivalent within $\pm0.01$; unmarked: inconclusive.}
\label{tab:scorer-comparison}
\resizebox{\columnwidth}{!}{\begin{tabular}{lrrrrrr}
\toprule
$-$ SBERT & \multicolumn{2}{c}{MovieLens} & \multicolumn{2}{c}{Amazon} & \multicolumn{2}{c}{Yelp} \\
\cmidrule(lr){2-3}\cmidrule(lr){4-5}\cmidrule(lr){6-7}
 & LGBM & MLP & LGBM & MLP & LGBM & MLP \\
\midrule
TF-IDF & -0.014$^{*}$ & -0.026$^{*}$ & -0.005 & -0.043$^{*}$ & -0.013$^{*}$ & -0.015$^{*}$ \\
mxbai (335M) & -0.007 & -0.008 & +0.003 & -0.015 & +0.003$^{\equiv}$ & +0.010$^{*}$ \\
phi4-mini raw (3.8B) & -0.104$^{*}$ & -0.135$^{*}$ & -0.015 & -0.101$^{*}$ & +0.002$^{\equiv}$ & -0.037$^{*}$ \\
Qwen3-Emb.\ (4B) & -0.002 & -0.012 & +0.010 & -0.002 & +0.006$^{*}$ & +0.016$^{*}$ \\
Hybrid & -0.010$^{*}$ & -0.023$^{*}$ & -0.003 & -0.066$^{*}$ & +0.007$^{*}$ & -0.020$^{*}$ \\
\bottomrule
\end{tabular}
}
\end{table}

Table~\ref{tab:scorer-comparison} reads the personalized,
popularity-matched configuration as a ladder of capacity steps. On
MovieLens (first two columns), moving
from TF-IDF to SBERT (22M parameters) is a significant gain
($+0.014$, $p<0.001$). It is the one step that pays: TF-IDF only ties
PureSVD, while SBERT beats it ($+0.017$, $p=0.011$). No step above SBERT
produces a detectable gain. mxbai (15$\times$ larger) and
Qwen3-Embedding-4B (180$\times$ larger) land within 0.007 of SBERT in
point estimate with CIs straddling zero. With 610 users, however, TOST
cannot establish equivalence at $\pm0.01$ either. We therefore claim
\emph{no detectable gain} above SBERT, not equivalence.

\textbf{A raw LLM is worse than no LLM.} Mean-pooled phi4-mini is the
outlier. It is $0.104$ below SBERT ($p<10^{-60}$), below TF-IDF, and
below PureSVD, which uses no text at all. Qwen3-Embedding-4B has
nearly the same size and the same quantization, and it recovers fully.
The deficit therefore comes from using an untrained chat model's hidden
states as an embedding, not from LLM capacity. Concatenating TF-IDF to
phi4-mini (Hybrid) recovers most of the loss but stays significantly
below SBERT alone ($-0.010$, $p=0.037$).

\textbf{Two scorers, one ordering.} The pointwise MLP reproduces this
ordering: SBERT is above TF-IDF ($+0.026$, $p<10^{-4}$), mxbai and
Qwen3 are not detectably different from SBERT, and phi4-mini ($-0.135$)
and Hybrid ($-0.023$) are significantly worse. The scorers disagree
about absolute levels, as expected for a listwise versus a pointwise
objective, but not about which representations help.

\subsection{Review text: Amazon and Yelp}
\label{sec:amazon-results}
\begin{table*}[t]
\centering
\caption{Amazon Reviews and Yelp NDCG@10 (LGBMRanker, mean $\pm$ 95\% CI
over 5 seeds). ``Item only'' and ``+Aff.+Pop.'' use popularity-matched
negatives, as in Table~\ref{tab:main-results}; ``Uniform'' uses uniform
negatives.}
\label{tab:amazon-results}
\begin{tabular}{lcccccc}
\toprule
 & \multicolumn{3}{c}{Amazon (249 users)} & \multicolumn{3}{c}{Yelp (10,537 users)} \\
\cmidrule(lr){2-4}\cmidrule(lr){5-7}
Pipeline & Item only & +Aff.+Pop. & Uniform & Item only & +Aff.+Pop. & Uniform \\
\midrule
Popularity & 0.090$\pm$0.003 & 0.090$\pm$0.003 & 0.132$\pm$0.012 & 0.133$\pm$0.001 & 0.133$\pm$0.001 & 0.465$\pm$0.003 \\
PureSVD & 0.151$\pm$0.007 & 0.151$\pm$0.007 & 0.166$\pm$0.007 & 0.151$\pm$0.003 & 0.151$\pm$0.003 & 0.298$\pm$0.002 \\
\midrule
TF-IDF & 0.068$\pm$0.009 & 0.170$\pm$0.011 & 0.180$\pm$0.005 & 0.110$\pm$0.003 & 0.179$\pm$0.001 & 0.465$\pm$0.003 \\
SBERT (22M) & 0.066$\pm$0.008 & 0.175$\pm$0.013 & 0.195$\pm$0.010 & 0.110$\pm$0.003 & 0.192$\pm$0.002 & 0.469$\pm$0.001 \\
mxbai (335M) & 0.064$\pm$0.009 & 0.178$\pm$0.015 & 0.187$\pm$0.010 & 0.109$\pm$0.002 & 0.195$\pm$0.001 & 0.472$\pm$0.002 \\
phi4-mini raw (3.8B) & 0.065$\pm$0.009 & 0.160$\pm$0.008 & 0.159$\pm$0.010 & 0.111$\pm$0.002 & 0.194$\pm$0.001 & 0.465$\pm$0.002 \\
Qwen3-Emb.\ (4B) & 0.065$\pm$0.010 & 0.185$\pm$0.007 & 0.196$\pm$0.013 & 0.110$\pm$0.002 & 0.198$\pm$0.002 & 0.469$\pm$0.003 \\
Hybrid & 0.066$\pm$0.009 & 0.172$\pm$0.008 & 0.186$\pm$0.007 & 0.110$\pm$0.002 & 0.199$\pm$0.001 & 0.470$\pm$0.003 \\
\bottomrule
\end{tabular}

\end{table*}

Table~\ref{tab:amazon-results} reports the two review datasets. Amazon
repeats the MovieLens pattern in direction but has little power. Every
personalized content pipeline beats popularity ($p<10^{-4}$), but with
249 evaluated users every LGBMRanker comparison between content
pipelines is inconclusive. The MLP scorer resolves only the extremes:
SBERT beats TF-IDF ($+0.043$) and phi4-mini is $0.101$ below SBERT.

Yelp, with 10{,}537 users, has enough power to resolve differences well
under 0.01. The first capacity step replicates: SBERT beats TF-IDF by
$0.013$ under LGBMRanker and $0.015$ under the MLP (both $p<10^{-12}$).
Unlike on the other datasets, capacity above SBERT also pays, but by
less. Qwen3-Embedding-4B is $+0.006$ above SBERT under LGBMRanker
($p=0.002$; 0.198 vs.\ 0.192, $+3\%$ relative) and $+0.016$ under the MLP
($+9\%$). mxbai is TOST-equivalent to SBERT under LGBMRanker and
$+0.010$ under the MLP. Raw phi4-mini is equivalent to SBERT under
LGBMRanker but $0.037$ below it under the MLP; \S\ref{sec:failure-analysis}
examines why. Across the three datasets, then, the TF-IDF-to-SBERT step
is the most consistent gain (on Amazon it is significant only under the
MLP). A trained LLM embedder adds a
further gain, half the size of that step under LGBMRanker and the same
size under the MLP, but only on long review text with a sample large
enough to detect it.

\subsection{Text redundancy, within one dataset}
\label{sec:redundancy-analysis}
\label{fig:redundancy-vs-gap}
Corpus lexical redundancy (mean pairwise TF-IDF cosine between items)
is a cheap candidate predictor of when dense encoders beat TF-IDF; in
retrieval, dense models' advantage over BM25 shrinks on high-overlap
corpora \cite{rare-redundancy}. Comparing datasets cannot isolate it,
because datasets also differ in domain, size, and text type, so we test
it within MovieLens. We
hold users, items, labels, splits, and negatives fixed and vary only
the item text:

\begin{center}
\small
\begin{tabular}{lccc}
\toprule
Item text & Redundancy & SBERT $-$ TF-IDF & $p$ \\
\midrule
Title only & 0.009 & $+0.019\pm0.008$ & $2\times10^{-6}$ \\
Title + genres & 0.030 & $+0.014\pm0.007$ & $5\times10^{-5}$ \\
Genres only & 0.188 & $-0.008\pm0.004$ & $3\times10^{-4}$ \\
\bottomrule
\end{tabular}
\end{center}

\noindent The dense encoder's advantage shrinks as redundancy rises and
reverses at the redundant extreme. All three differences survive a
Bonferroni correction. The manipulation is not a pure redundancy
intervention: genres-only text also removes the title, and with
$\sim$20 genre tokens TF-IDF becomes an exact genre indicator. The
result therefore shows that the lexical-versus-dense ordering
\emph{follows} text redundancy within a dataset, which is the property
a practitioner can measure before running anything. It does not show
that redundancy alone causes the ordering. Nor does it transfer across
datasets: Yelp's review text is the most redundant corpus we test
(0.236; Amazon 0.193), yet SBERT beats TF-IDF there by 0.013. Redundancy
is a guide for choosing which text to embed within one corpus, not a
predictor that ranks corpora.

\subsection{Robustness under temporal shift}
\label{sec:robustness}
\label{tab:robustness}
\label{fig:robustness}
Holding out each user's most recent positives lowers every pipeline
(popularity-matched negatives, LGBMRanker, 5 seeds). On MovieLens, NDCG@10
drops by 0.115--0.128, and the ordering established under the random
split holds. TF-IDF degrades most (0.576 $\rightarrow$ 0.448, a drop of 0.128),
while SBERT and mxbai degrade least (0.115 each; SBERT 0.589
$\rightarrow$ 0.474), so the SBERT-over-TF-IDF gap roughly doubles, from $+0.014$ to $+0.027$ ($p<10^{-4}$). mxbai also separates
from TF-IDF under shift ($+0.020$, $p=0.006$). No encoder above SBERT is
detectably better than it (mxbai $-0.006$, Qwen3 $-0.010$, both
$p>0.5$), and raw phi4-mini remains far behind ($-0.115$). The only
reliable robustness gain therefore again comes from the first step
off lexical features. On Amazon, every pipeline drops by a similar
amount (0.055--0.069), every content pipeline still beats popularity
($p<10^{-3}$), and no encoder is measurably more robust than another.
On Yelp the shift costs content pipelines only 0.005--0.009 (popularity
0.016). Both gains survive it: SBERT over TF-IDF ($+0.011$) and Qwen3
over SBERT ($+0.010$, $p<10^{-4}$). Among content pipelines, Qwen3 loses the least (0.005).

\subsection{Quality--cost}
\label{sec:quality-cost}

Table~\ref{tab:latency} reports per-item embedding latency from the
like-for-like benchmark (\S\ref{sec:experimental-setup}).

\begin{table}[t]
\centering
\caption{Median per-item embedding latency (ms) over three rounds, 256
items, full-text coverage. Ratios are relative to SBERT at the same
batch size. MovieLens text is a title and genres; Amazon text averages
about 2{,}100 characters of review snippets.}
\label{tab:latency}
\footnotesize
\setlength{\tabcolsep}{2pt}
\begin{tabular}{lrrrr}
\toprule
 & \multicolumn{2}{c}{MovieLens} & \multicolumn{2}{c}{Amazon} \\
\cmidrule(lr){2-3}\cmidrule(lr){4-5}
Encoder & batch 1 & batch 32 & batch 1 & batch 32 \\
\midrule
SBERT (22M) & 12.3 & 5.8 & 60 & 42 \\
mxbai (335M) & 42.5 (3.5$\times$) & 28.8 (5.0$\times$) & 257 (4.3$\times$) & 241 (5.7$\times$) \\
phi4-mini (3.8B) & 222 (18$\times$) & 170 (29$\times$) & 2{,}691 (45$\times$) & 2{,}696 (64$\times$) \\
Qwen3-Emb.\ (4B) & 231 (19$\times$) & 262 (45$\times$) & 2{,}974 (49$\times$) & 2{,}924 (70$\times$) \\
\bottomrule
\end{tabular}
\end{table}

With the same inputs, batching, and hardware for every encoder, the
two 4B-class models cost 18--45$\times$ SBERT per item on short
MovieLens text and 45--70$\times$ on long Amazon text, and mxbai costs
3.5--5.7$\times$. The gap widens with text length because SBERT and
mxbai process short chunks while the LLMs attend over the full text,
and widens with batch size because SBERT gains more from batching
than the LLM servers do. On Amazon, batching barely helps the LLMs:
with four parallel server slots on one machine, long sequences
saturate the accelerator even at batch size 1. Qwen3 at batch 32 on
MovieLens is slower per item than at batch 1, which we attribute to
contention among the server slots, a property of this deployment
rather than of the model. A naive comparison of SBERT's batched corpus
encoding against unbatched, sequential LLM calls inflates the ratio by
more than an order of magnitude, so latency claims should state the
batching regime.

Set against quality, the extra cost buys nothing measurable. On
MovieLens, Qwen3 costs 19--45$\times$ SBERT for an NDCG@10 difference of
$-0.002$ (inconclusive), mxbai costs 3.5--5$\times$ for $-0.006$, and
phi4-mini costs 18--29$\times$ for $-0.104$. On Amazon, Qwen3 has the
highest point estimate ($+0.010$ over SBERT under LGBMRanker, $-0.002$
under the MLP), but neither difference is detectable at 45--70$\times$
the cost. Yelp is the one case where capacity above SBERT buys a
detectable gain: $+0.006$ to $+0.016$ NDCG@10 for Qwen3. Yelp's item text
is as long as Amazon's (mean 2{,}148 vs.\ 2{,}083 characters), so its
per-item cost should be similar, about 45--70$\times$ SBERT. Because item embeddings are computed once per item and
cached, these ratios bound refresh and cold-start cost rather than
request latency (\S\ref{sec:practical-guidelines}).

\section{Where Capacity Helps and Where It Fails}
\label{sec:failure-analysis}

\subsection{When raw LLM hidden states fail}
Mean-pooled phi4-mini is the least predictable encoder: far below
TF-IDF on MovieLens, but level with SBERT on Yelp under LGBMRanker
(Table~\ref{tab:scorer-comparison}). For each embedding space we
measured the effective dimensionality (participation ratio of the
covariance spectrum) on up to 1{,}000 random items. We also measured
how well profile similarity \emph{alone} ranks each user's test
candidates: candidates are sorted by their mean cosine to the user's
training positives, with no scorer (seed 1, popularity-matched
negatives). The second measure depends only on the order of the
similarities, not on their range.

\begin{center}
\small
\setlength{\tabcolsep}{3pt}
\begin{tabular}{lrrrcrrr}
\toprule
 & \multicolumn{3}{c}{Effective dim.} & & \multicolumn{3}{c}{Affinity-only NDCG@10} \\
\cmidrule(lr){2-4}\cmidrule(lr){6-8}
Space & ML & Amz & Yelp & & ML & Amz & Yelp \\
\midrule
TF-IDF & 198 & 125 & 207 & & 0.393 & 0.193 & 0.140 \\
SBERT & 68 & 28 & 37 & & 0.410 & 0.183 & 0.152 \\
mxbai & 57 & 40 & 42 & & 0.419 & 0.178 & 0.160 \\
phi4-mini (raw) & \textbf{21} & \textbf{3.7} & \textbf{3.6} & & \textbf{0.198} & 0.168 & 0.157 \\
Qwen3-Emb. & 92 & 51 & 95 & & 0.411 & 0.204 & 0.168 \\
\midrule
Popularity & & & & & 0.379 & 0.090 & 0.133 \\
\bottomrule
\end{tabular}
\end{center}

\noindent Raw phi4-mini vectors are nearly parallel on every dataset
(mean pairwise cosine 0.95--0.96), which is the anisotropy that
contrastive training removes \cite{sgpt, llm2vec}. Qwen3-Embedding has
that training, and its spread is comparable to the dedicated encoders'.
Geometry alone does not decide the outcome, however. On short MovieLens
metadata, the \emph{order} of phi4-mini's similarities is close to
uninformative: ranking by it scores 0.198, below popularity, while every
other encoder scores 0.39--0.42. Mean-centering does not repair it
(0.481 $\rightarrow$ 0.486 NDCG@10 in the full pipeline, against SBERT's
0.574 on the same seed). On long Yelp reviews the order is as
informative as mxbai's and better than SBERT's. LGBMRanker splits on
feature order, so it inherits this pattern: phi4-mini fails on
MovieLens and matches SBERT on Yelp. The MLP scorer fails with
phi4-mini on all three datasets. One explanation consistent with this,
which we did not test directly, is that the MLP standardizes every
input dimension. That gives thousands of near-constant item-vector
dimensions the same scale as the few informative ones, while a tree
model can ignore them. In short, whether raw hidden states work depends
on text length and on the scorer. An LLM of the same size trained for
embedding has neither dependence.

\subsection{User history does not moderate the gain}
On MovieLens, an exploratory per-user regression suggested that the
advantage of mxbai and Qwen3 over SBERT grows with the log of a user's
number of positives (slope $+0.004$ to $+0.019$ per unit, $p<0.05$ in
all twelve configuration--scorer--encoder combinations). The pattern
did not replicate. On Amazon all slopes have $p>0.17$. On Yelp the
LGBMRanker slopes are positive but not significant ($p\geq0.07$), and
the MLP slopes \emph{reverse} sign ($p<0.01$ in five of six). We
therefore do not treat history length as a moderator of encoder
capacity.

\section{Practical Guidelines}
\label{sec:practical-guidelines}

The findings above give an ordered checklist for deciding whether to
put a larger encoder, or an LLM, into an existing ranking pipeline.

\begin{enumerate}
  \setcounter{enumi}{-1}
  \item \textbf{Check the evaluation before the encoder.} Make sure the
    scorer actually sees the user, for example through
    profile-similarity features, and evaluate with negatives that
    popularity alone cannot separate. Report a popularity baseline and a
    collaborative baseline. Without these three things every encoder
    can look the same, and the comparison is uninformative
    (\S\ref{sec:movielens-results}).
  \item \textbf{Choose which text to embed.} Within one corpus,
    fields with highly redundant text (e.g.\ a small tag vocabulary)
    favored TF-IDF, and distinctive fields favored dense encoders.
    Mean pairwise TF-IDF cosine measures this without training. It does
    not compare corpora: Yelp's highly redundant reviews still favored
    the dense encoder (\S\ref{sec:redundancy-analysis}).
  \item \textbf{Try the smallest dense encoder first.} The step from
    TF-IDF to a 22M-parameter encoder was the most consistent gain,
    significant on all three datasets. Benchmark every larger encoder against that small one,
    not against TF-IDF. Otherwise the small encoder's gain gets credited
    to the large one.
  \item \textbf{Never use a chat LLM's raw hidden states as an
    embedding.} Mean-pooled phi4-mini was worse than TF-IDF and than a
    text-free collaborative model on short metadata, and failed under
    the MLP scorer everywhere. Centering did not fix it. If an LLM-scale
    encoder is wanted, use one trained for embedding
    (\S\ref{sec:failure-analysis}).
  \item \textbf{Size the test to the gain you are looking for.} Gains
    above the small encoder were at most 0.016 NDCG@10 and needed
    $\sim$10K users to detect. With a few hundred users, report an
    equivalence test, and treat a non-significant result as no
    evidence, not as no gain (\S\ref{sec:amazon-results}).
  \item \textbf{Price the serving path.} Item embeddings are content
    features, so they are computed once per item update and cached, and
    a large per-item latency matters little offline. It becomes a hard
    constraint only when items change often or must be embedded at
    request time. Measured like-for-like, a 4B model cost 18--70$\times$
    SBERT per item here, more on long text, for at most 0.016 NDCG@10
    (\S\ref{sec:quality-cost}).
  \item \textbf{Re-check under temporal shift.} Use a split that holds
    out each user's most recent interactions, not only a random one
    (\S\ref{sec:robustness}).
\end{enumerate}

\section{Discussion}
\label{sec:discussion}

\subsection{What the results say about capacity}
Treating ``add an LLM'' as a hypothesis, not a default upgrade, holds up
under a stricter evaluation, but for different reasons than a naive
comparison would suggest. Most of the value of a text encoder in a
personalized ranker comes from having \emph{some} dense, well-spread
representation in which a user's history can be compared to a
candidate. A 22M-parameter encoder already provides most of it.
Additional capacity paid off only on long review text, by a few
thousandths to a hundredth of NDCG@10, and only from an encoder trained
to embed. Raw LLM hidden states are a poor substitute: their
near-parallel geometry makes profile similarity fragile, uninformative
on short text and unusable by a scorer that standardizes its inputs.

These results complement rather than contradict work that reports gains
from language-model item representations
\cite{unisrec, recformer, yuan-id-vs-modality, harte-llmseqrec,
amazon-reviews-2023}. Those methods train or fine-tune the encoder
jointly with the recommender, or pretrain it on in-domain interaction
data. Our encoders are frozen and generic, which is the cheaper and
more common deployment path, and the one for which our checklist is
written.

\subsection{Limitations}
\label{sec:limitations}
\textbf{Scale and datasets.} Amazon's 249 evaluated users are too few
to resolve most pairwise differences, and MovieLens's 610 cannot
establish equivalence among SBERT, mxbai, and Qwen3 at $\pm0.01$. Only
Yelp has the power to detect gains below 0.01, so the gain above SBERT
rests on one dataset, from one city, with short user histories (median
3 positives). It may be larger or smaller in domains with longer
histories.

\textbf{Models.} Each rung of the ladder is one model: one small
encoder, one large encoder, one raw chat LLM, and one trained LLM
embedder, all run locally and quantized. We do not evaluate hosted
embedding APIs, other model families or sizes, or fine-tuning any
encoder on the target domain.

\textbf{Evaluation protocol.} Metrics are computed over sampled
candidate lists, not full-catalog ranking. Popularity-matched negatives
reduce, but do not remove, the leniency of sampled metrics
\cite{krichene-sampled}. Scorer hyperparameters are fixed across
pipelines rather than tuned per pipeline. This keeps the comparison
controlled but could disadvantage representations that need different
settings. Under the temporal split, popularity features can include
other users' later ratings (\S\ref{sec:candidate-generation}).

\textbf{Analyses.} The redundancy manipulation changes information
content along with redundancy, and it does not transfer across datasets
(\S\ref{sec:redundancy-analysis}). Our account of why the MLP fails
with raw phi4-mini is consistent with the data but untested
(\S\ref{sec:failure-analysis}).

\section{Conclusion}
\label{sec:conclusion}

We asked when spending more model capacity on item text pays off inside
a learning-to-rank pipeline. We compared six frozen encoders, from
TF-IDF to a 4B LLM trained for embedding, on three datasets under a
personalized scorer, hard and easy negatives, random and temporal
splits, two scorers, paired equivalence testing, and a like-for-like
latency benchmark.

Three conclusions follow. First, evaluation design can erase the
question: a scorer blind to the user collapses to popularity, and easy
negatives let popularity match every encoder. Second, once those are
fixed, the cheapest dense encoder captures the most consistent gain.
Larger encoders add little: nothing detectable on MovieLens or Amazon,
and 0.006--0.016 NDCG@10 on Yelp at roughly 45--70$\times$ the per-item
latency. Third, a chat LLM's raw hidden states are not a safe embedding.
They can fall below TF-IDF, depending on text length and scorer, a
failure that an LLM of the same size trained for embedding does not
share.

\textbf{Future work.} Online evaluation through interleaving,
replication in domains with long user histories, and fine-tuning the
small encoder on in-domain interactions, which may close the remaining
gap at a fraction of the cost.

\section*{AI Disclosure}
We used Claude (Anthropic), to assist with the paper and code. The
authors take full responsibility for the accuracy, originality, and
integrity of the submitted work.


\bibliographystyle{IEEEtran}
\bibliography{references}

\end{document}
